%% kapa: exam=1 round=19 year=2008 period=1 form=A subject=appraisal_law subject_ko=부동산관계법규 seq=20 number=60
%% source: Q-Net 감정평가사 기출문제 (cst00309) artlSeq=1001962 file=extracted/1008/1001962/2008년 제19회 감정평가사 1차시험 기출문제/1교시.hwp
%% conversion: kordoc 4.24.51 -> tools/gen1.py (auto); review=proofread-sonnet; flags=none
\begin{question}{60}
감정평가업자의 업무와 관련된 판례의 태도로 옳지 않은 것은?
\begin{choices}
  \item 감정평가업자가 현장조사 당시 감정대상 주택소유자의 처로부터 임대차가 없다는 확인을 받고 감정평가서에 ‘임대차 없음’이라고 기재하였으나 이후에 임차인의 존재가 밝혀진 경우 감정평가업자는 손해배상책임을 부담하지 않는다.
  \item 감정평가업자가 과실로 감정평가 당시의 적정가격과 현저한 차이가 있게 감정평가함으로써 감정평가의뢰인에게 손해를 발생하게 한 때에는 그 손해를 배상할 책임이 있다.
  \item 감정평가업자가 부실감정으로 인하여 감정평가의뢰인이나 선의의 제3자에게 손해를 입힌 경우 (구)지가공시 및 토지 등의 평가에 관한 법률상 손해배상책임과 민법상의 불법행위로 인한 손해배상책임을 함께 물을 수 있다.
  \item 감정평가업자가 금융기관으로부터 조사를 의뢰받은 담보물건과 관련된 임대차관계 등을 조사함에 있어 단순히 다른 조사기관의 전화조사만으로 확인된 내용을 기재한 임대차확인조사서를 제출한 것은 감정평가업자에게 감정평가업무협약에 따른 조사의무를 다하지 아니한 과실이 있다.
  \item 감정평가업자가 실지조사․공부조사 등 감정평가에 수반되는 조사업무를 통상적으로 요구되는 주의정도에 따라 성실히 수행하였음에도 당해 토지가 하천구역에 편입되어 국유로 된 토지인 사실을 알아내지 못한 채 그 시가에 대하여 감정평가를 하였다고 하더라도 이를 가지고 과실로 인한 감정평가의 하자라고 볼 수는 없다.
\end{choices}
\end{question}
\answer{1}
